% Gemini theme --- https://github.com/anishathalye/gemini
% Poster for ICLR 2026

\documentclass[final]{beamer}

% ====================
% Packages
% ====================

\usepackage[T1]{fontenc}
\usepackage[size=custom,width=194,height=95,scale=1.55]{beamerposter}
\usetheme{gemini}
\usecolortheme{gemini}
\usepackage{graphicx}
\usepackage{booktabs}
\usepackage{tikz}
\usepackage{amsmath,amssymb,mathtools}
\usepackage{newtxmath}
% Fix mathrm/operatorname to use Times (fontspec overrides it to Lato)
\SetSymbolFont{operators}{normal}{OT1}{ntxtlf}{m}{n}
\SetSymbolFont{operators}{bold}{OT1}{ntxtlf}{b}{n}
\usepackage{xcolor}
\usepackage{tcolorbox}
\usepackage{qrcode}
\tcbuselibrary{skins}

% ====================
% Custom Colors (project branding)
% ====================

% Color palette
\definecolor{SeaGreen}{HTML}{008f69}
\definecolor{BattleshipGray}{HTML}{848482}
\definecolor{AirBlue}{HTML}{71a0c1}
\definecolor{HonoluluBlue}{HTML}{0e79b2}
\definecolor{NavyBlue}{HTML}{000080}
\definecolor{Charcoal}{HTML}{2c3e50}
\definecolor{LightLavender}{HTML}{f0f4ff}

% Override Gemini colors — neutral header, accent blue details
\setbeamercolor{headline}{fg=white,bg=HonoluluBlue}
\setbeamercolor{headline title}{fg=white}
\setbeamercolor{headline author}{fg=white}
\setbeamercolor{headline institute}{fg=white}
\setbeamercolor{headline rule}{bg=HonoluluBlue}
\setbeamercolor{block title}{fg=NavyBlue,bg=white}
\setbeamercolor{block separator}{bg=AirBlue}
\setbeamercolor{block body}{fg=Charcoal,bg=white}
\setbeamercolor{block title alerted}{fg=white,bg=HonoluluBlue}
\setbeamercolor{block separator alerted}{bg=AirBlue}
\setbeamercolor{block body alerted}{fg=Charcoal,bg=LightLavender}
\setbeamercolor{item}{fg=HonoluluBlue}
\setbeamercolor{heading}{fg=HonoluluBlue}
\setbeamercolor{footline}{fg=white,bg=HonoluluBlue}

% ====================
% Callout box style
% ====================

\newtcolorbox{resultbox}[1][]{
  colback=AirBlue!15,
  colframe=AirBlue!15,
  boxrule=0pt,
  arc=3pt,
  left=6pt, right=6pt, top=4pt, bottom=4pt,
  #1
}

% ====================
% Lengths
% ====================

\newlength{\sepwidth}
\newlength{\colwidth}
\setlength{\sepwidth}{0.02\paperwidth}
\setlength{\colwidth}{0.3067\paperwidth}

% ====================
% Title
% ====================

\title{Shoot First, Ask Questions Later? Building Rational Agents that Explore and Act Like People}

\author{Gabriel Grand\textsuperscript{*1,2} \quad Valerio Pepe\textsuperscript{*3} \quad Joshua B.\ Tenenbaum\textsuperscript{1,2} \quad Jacob Andreas\textsuperscript{1}}

\institute{\textsuperscript{1}MIT CSAIL \quad \textsuperscript{2}MIT Brain and Cognitive Sciences \quad \textsuperscript{3}Harvard SEAS \qquad {\normalsize\textsuperscript{*}Equal contribution}}

% \logoleft{\includegraphics[height=7cm]{logos/mit-csail.png}}
% \logoright{\includegraphics[height=7cm]{logos/qr-code.png}}

% ====================
% Footer
% ====================

\footercontent{
  \textbf{gabegrand.github.io/battleship} \hfill
  ICLR 2026 --- Oral Presentation \hfill
  github.com/gabegrand/battleship
}

% ====================
% Document
% ====================

\begin{document}
\begin{frame}[t]
\begin{columns}[t]

% ====================
% Column 1: Motivation + Task + BED Framework
% ====================

\begin{column}{\sepwidth}\end{column}
\begin{column}{\colwidth}

\begin{block}{Motivation}

In settings like \textbf{scientific discovery} and \textbf{medical diagnosis}, AI agents must \emph{ask smart questions} before acting---but current LMs are optimized to \emph{answer} questions, not \emph{ask} good ones.

\end{block}

\begin{block}{Collaborative Battleship}

\begin{center}
  \includegraphics[width=0.95\colwidth]{figures/battleship-splash-2.png}
\end{center}

\begin{itemize}
  \item \textbf{Captain}: asks yes/no questions, then fires at tiles
  \item \textbf{Spotter}: sees full board, provides grounded answers
  \item 8$\times$8 board, 4 ships, budget of 15 questions + 40 shots
\end{itemize}

\end{block}

\begin{block}{Bayesian Experimental Design}

Given hidden state $S \in \mathcal{S}$, experiment $x \in X$, observation $y$:
%
\begin{equation*}
  \underbrace{p(S \mid x, y)}_{\text{posterior}}
  \;\propto\;
  \underbrace{p(y \mid S, x)}_{\text{likelihood}}
  \;\;
  \underbrace{p(S)}_{\text{prior}}
\end{equation*}

\heading{Expected Information Gain}

\begin{align*}
  \operatorname{IG}(x, y) &= H\big[p(S)\big] - H\big[p(S \mid x, y)\big] \\[0.5ex]
  \operatorname{EIG}(x) &= \mathbb{E}_{y \sim p(y \mid x)}\!\big[\operatorname{IG}(x, y)\big], \quad
  x^{\star} = \arg\max_{x \in X}\; \operatorname{EIG}(x)
\end{align*}

For binary answers with a noisy oracle ($\mathrm{BSC}(\varepsilon)$):
%
\begin{equation*}
  \boxed{\;\operatorname{EIG}_\varepsilon(S;\, q) = H_b\!\big(\varepsilon + (1 - 2\varepsilon)\,p\big) - H_b(\varepsilon)\;}
\end{equation*}

where $p := P(a\!=\!1) = \mathbb{E}_{s \sim P(S)}\!\big[P(a \mid q;\, s)\big]$.
Maximized when $p \approx 0.5$.

\end{block}

\end{column}

% ====================
% Column 2: Strategies + SpotterQA + CaptainQA
% ====================

\begin{column}{\sepwidth}\end{column}
\begin{column}{\colwidth}

\begin{block}{Bayesian Strategies}

Three inference-time strategies leveraging the BED framework:

\heading{$Q_{\mathrm{Bayes}}$: Maximize information gain}

Sample $K$ candidates from the LM, select the most informative:
%
\begin{equation*}
  q_t^{\star} \in \arg\max_{q \in \mathcal{Q}}\; \operatorname{EIG}_\varepsilon(q \mid b_t)
\end{equation*}

\heading{$M_{\mathrm{Bayes}}$: Maximize hit probability}

Select the tile most likely to contain a ship:
%
\begin{equation*}
  u_t^{\star} \in \arg\max_{u}\; p_{\mathrm{hit}}(u \mid b_t), \;\; p_{\mathrm{hit}}(u \mid b_t) \coloneqq \!\!\sum_{s \in \mathcal{F}} b_t(s)\, \mathbb{1}[u \text{ is ship in } s]
\end{equation*}

\heading{$D_{\mathrm{Bayes}}$: Explore vs.\ exploit}

One-step lookahead---ask if expected post-question hit prob.\ exceeds current:
%
\begin{equation*}
  \text{Ask } q_t^{\star} \;\text{if}\; \gamma\,\widehat{p}_{\mathrm{hit}}^{\,t+1}(q_t^{\star}) \;>\; p_{\mathrm{hit}}(u_t^{\star} \mid b_t)
\end{equation*}

\end{block}

\begin{block}{SpotterQA: Question Answering}

931 gold-labeled questions, \textbf{15 LMs}, 4 strategies.
Bars = mean accuracy $\pm$ s.d.\ across models. Dashed line = human.

\begin{center}
  \includegraphics[height=16cm]{figures/spotter_accuracy_overall.pdf}
\end{center}

\begin{resultbox}
  \textbf{Code generation boosts accuracy by +14.7\%} over LM baselines.
  Best model reaches 92.8\% (human: 92.5\%).
\end{resultbox}

\end{block}

\end{column}

% ====================
% Column 3: CaptainQA + Generalization + References
% ====================

\begin{column}{\sepwidth}\end{column}
\begin{column}{\colwidth}

\begin{block}{CaptainQA: Full Game Performance}

54 games, 3 LMs + strategy ablations. \enspace
{\color{AirBlue}$\blacksquare$}\,Llama-4-Scout \enspace {\color{HonoluluBlue}$\blacksquare$}\,GPT-4o \enspace {\color{NavyBlue}$\blacksquare$}\,GPT-5

\begin{columns}[T]
  \begin{column}{0.48\colwidth}
    \begin{center}
      \includegraphics[height=14cm]{figures/captain_f1_boxplot.pdf}
    \end{center}
  \end{column}
  \begin{column}{0.48\colwidth}
    \begin{center}
      \includegraphics[height=14cm]{figures/eig_max_vs_k.pdf}
    \end{center}
  \end{column}
\end{columns}

\begin{resultbox}
  \begin{center}
  \begin{tabular}{ccc}
    \textbf{Llama-4-Scout F1} & \textbf{vs.\ Humans} & \textbf{vs.\ GPT-5} \\[0.3ex]
    {\LARGE\color{HonoluluBlue}\bfseries 0.37 $\rightarrow$ 0.76} &
    {\LARGE\color{SeaGreen}\bfseries 82\% win} &
    {\LARGE\color{SeaGreen}\bfseries 67\% win} \\
  \end{tabular}
  \end{center}
  \vspace{-0.5ex}
  {\small \textbf{99.7\% cost savings} vs.\ GPT-5 (\$8 vs.\ \$874). EIG: \textbf{+0.227 bits} (94\% of ceiling) with $K$=10. Redundant questions: \textbf{18.5\% $\rightarrow$ 0.2\%}.}
\end{resultbox}

\end{block}

\begin{block}{Generalization: Guess Who?}

Same strategies on \textbf{Guess Who?} (100 characters, 8 questions).
Llama: 0.300 $\rightarrow$ \textbf{0.724}; \enspace GPT-4o: 0.617 $\rightarrow$ \textbf{0.900}.

\begin{center}
  \includegraphics[height=8cm]{figures/guess-who-results.pdf}
\end{center}

\end{block}

\begin{block}{References}

{%
\footnotesize
\begin{minipage}[c]{0.82\colwidth}
\begin{tabular}{@{}r@{\;\;}p{0.92\linewidth}@{}}
{[1]} & A.~Rothe, B.M.~Lake, T.M.~Gureckis. Question asking as program generation. \textit{NeurIPS}, 2017. \\
{[2]} & A.~Rothe, B.M.~Lake, T.M.~Gureckis. Do people ask good questions? \textit{CBB}, 2018. \\
{[3]} & A.~Rothe, B.M.~Lake, T.M.~Gureckis. Asking goal-oriented questions and learning from answers. \textit{CogSci}, 2019. \\
{[4]} & G.~Grand, V.~Pepe, J.~Andreas, J.B.~Tenenbaum. Loose LIPS sink ships. \textit{CogSci}, 2024. \\
{[5]} & C.E.~Shannon. A mathematical theory of communication. \textit{Bell Syst.\ Tech.\ J.}, 1948. \\
{[6]} & D.V.~Lindley. On a measure of the information provided by an experiment. \textit{Ann.\ Math.\ Stat.}, 1956. \\
{[7]} & T.~Rainforth, A.~Foster, D.R.~Ivanova, F.B.~Smith. Modern Bayesian experimental design. \textit{arXiv}, 2023. \\
\end{tabular}
\end{minipage}\hfill
\begin{minipage}[c]{0.11\colwidth}
  \centering
  \qrcode[height=0.9\linewidth]{https://gabegrand.github.io/battleship}
  \par\vspace{0.5ex}
  {\sffamily\scriptsize\textbf{Paper \& Code}}
\end{minipage}%
}

\end{block}

\end{column}
\begin{column}{\sepwidth}\end{column}

\end{columns}
\end{frame}
\end{document}
